\Article{19}
\textsc{Theorem.} \textit{With everything as above, a circle can be drawn with centre $C$,}
\SourceLeaf{030}
\textit{on whose circumference there are $2m$ points where $T=0$ and as many where $U=0$, arranged so that each of the latter lies between two of the former.}

Let the sum of all the coefficients $A,B$ etc. $K,L,M$, taken positively, be \mbox{$=S$}, and take $R$ to be both $>S\BareRoot{2}$ and $>1$\HistoricalNote{9}. I then say that what the theorem states necessarily holds for the circle drawn with radius $R$. For brevity, label the point of this circumference that is $\frac1m45$ degrees from its intersection with the left part of the axis, or for which $\varphi=\frac1m45^\circ$, by (1); similarly, the point that is $\frac3m45^\circ$ from that intersection, or for which $\varphi=\frac3m45^\circ$, by (3); then the point where $\varphi=\frac5m45^\circ$ by (5), and so on up to $(8m-1)$, which is $\frac{8m-1}m45$ degrees from that intersection if one always proceeds in the same direction (or $\frac1m45^\circ$ in the opposite direction). Thus there are altogether $4m$ points on the circumference at equal intervals. Between $(8m-1)$ and (1) there will be one point where $T=0$; likewise, one such point will lie between (3) and (5), between (7) and (9), between (11) and (13), etc., so that their number is $2m$. In the same way, one point where $U=0$ will lie between (1) and (3), between (5) and (7), and between (9) and (11), so that their number is also $=2m$. Finally, apart from these $4m$ points, there will be no others anywhere on the circumference where either $T$ or $U$ is $=0$.

\textit{Proof.} I. At the point (1), $m\varphi=45^\circ$, and thus
\[ T=R^{m-1}(R\BareRoot{\tfrac12}+A\sin(m-1)\varphi+\tfrac BR\sin(m-2)\varphi+\text{etc.}+\tfrac L{R^{m-2}}\sin\varphi)
\]
But the sum $A\sin(m-1)\varphi+\frac BR\sin(m-2)\varphi$ etc. certainly cannot be greater than $S$, and must therefore be less than $R\BareRoot{\frac12}$. It follows that at this point the value of $T$ is certainly positive. All the more, then, $T$ will have a positive value when $m\varphi$ lies between $45^\circ$ and $135^\circ$; that is, the value of $T$ will always be positive from the point (1) to (3). For the same reason, $T$ will be positive everywhere from the point (9) to (11), and in general from any point $(8k+1)$ to $(8k+3)$, where $k$ denotes any integer. Similarly, $T$ will be negative everywhere between (5) and (7), between (13) and (15), etc., and in general between $(8k+5)$ and $(8k+7)$; thus it cannot be $=0$ anywhere in any of these intervals. But since its value is positive at (3) and negative at (5), it must be $=0$ somewhere between (3) and (5); likewise somewhere between (7) and (9), between (11) and (13), etc.,
\SourceLeaf{031}
up to and including the interval between $(8m-1)$ and (1). Thus we have $T=0$ at altogether $2m$ points. Q. E. P.

II. That there are no other points with this property apart from these $2m$ points is seen as follows. Since there are none between (1) and (3), between (5) and (7), etc., there could be more such points only if at least two lay in some interval between (3) and (5), or between (7) and (9), etc. Then, however, $T$ would necessarily be a \textit{maximum} or \textit{minimum} somewhere in the same interval, and hence $\frac{\mathrm dT}{\mathrm d\varphi}=0$. But $\frac{\mathrm dT}{\mathrm d\varphi}=mR^{m-2}(R\cos m\varphi+\frac{m-1}mA\cos(m-1)\varphi+\text{etc.})$, and between (3) and (5), $\cos m\varphi$ is always negative and $>\BareRoot{\frac12}$. It is therefore easy to see that $\frac{\mathrm dT}{\mathrm d\varphi}$ is negative throughout this interval; in the same way, it is positive everywhere between (7) and (9), negative between (11) and (13), etc. Thus it cannot be 0 in any of these intervals, and the supposition cannot stand. Therefore etc. Q. E. S.

III. In exactly the same way, it is proved that $U$ has a negative value everywhere between (3) and (5), between (11) and (13), etc., and in general between $(8k+3)$ and $(8k+5)$; but a positive value between (7) and (9), between (15) and (17), etc., and in general between $(8k+7)$ and $(8k+9)$. It follows at once that $U=0$ must hold somewhere between (1) and (3), between (5) and (7), etc., that is, at $2m$ points. But $\frac{\mathrm dU}{\mathrm d\varphi}=0$ cannot hold in any of these intervals (which is easily proved in the same way as above). There will therefore be no more than those $2m$ points on the circumference of the circle where $U=0$. Q. E. T. and Q.

The part of the theorem stating that there are no more than $2m$ points where $T=0$, and no more than $2m$ where $U=0$, can also be proved from the fact that equations $T=0$, $U=0$ represent curves of degree $m$. For it is known from higher geometry that such curves cannot be cut in more than $2m$ points by a circle, as a curve of the second degree.

\Article{20} If another circle is drawn from the same centre with a radius greater than $R$ and divided in the same way, then on this circle too there will be one point between (3) and (5) where $T=0$, and likewise between (7) and (9), etc. It is easy to see that the less the radius of this circle differs from the radius $R$, the nearer such points lying between (3) and (5) on the two circumferences must be to each other. The same will hold if a circle is drawn with a radius somewhat smaller than $R$, but still greater than $S\BareRoot{2}$ and 1. From this it is readily understood that the circumference of the circle drawn with radius $R$
\SourceLeaf{032}
is actually \textit{cut} by some branch of the first line at the point between (3) and (5) where $T=0$; and the same holds for the remaining points where $T=0$. In the same way, it is clear that the circumference of this circle is cut by some branch of the second line at all $2m$ points where $U=0$. These conclusions may also be expressed as follows: if a circle of the appropriate size is drawn with centre $C$, $2m$ branches of the first line and as many branches of the second line will enter it, in such a way that each pair of neighbouring branches of the first line is separated by some branch of the second line. See Fig. 2, where the circle will now be of finite rather than infinite size. The numbers attached to the individual branches must not be confused with the numbers by which, for brevity, I designated certain boundaries on the circumference in this and the preceding article.

\Article{21} From this relative arrangement of the branches entering the circle, the necessary existence of an intersection of some branch of the first line with a branch of the second line inside the circle can now be deduced in so many different ways that I almost do not know which method to choose in preference to the others. The following seems the clearest. Let us label (Fig. 2) the point on the circumference where it is cut by the left part of the axis (which is itself one of the $2m$ branches of the first line) by 0; the next point, where a branch of the second line enters, by 1; the point next to this, where the second branch of the first line enters, by 2; and so on up to $4m-1$. Thus a branch of the second line enters the circle at every point marked with an even number, while a branch of the second line enters at every point marked with an odd number. Now it is known from higher geometry that every algebraic curve (or each individual part of any algebraic curve, if it happens to consist of several parts) either returns into itself or extends to infinity in both directions. Hence, if a branch of an algebraic curve enters a bounded region, it must necessarily leave that region again somewhere\HistoricalNote{10}. From this it is easily concluded that every point marked with an even number\SourceLeaf{033}
 (or, for brevity, \textit{every even point}) must be joined within the circle to another even point by a branch of the first line; similarly, every point marked with an odd number must be joined to another such point by a branch of the second line. Although this connection of pairs of points may vary greatly with the nature of the polynomial $X$, so that it cannot be determined in general, it is nevertheless easy to prove that \textit{whatever it may be, an intersection of the first line with the second line always results.}

\Article{22} It seems that this necessity can most conveniently be proved by contradiction. Suppose, then, that each pair of even points and each pair of odd points can be joined in such a way that no intersection of a branch of the first line with a branch of the second line results. Since the axis is part of the first line, the point 0 will clearly be joined to the point $2m$. The point 1 therefore cannot be joined to any point lying beyond the axis, that is, to any point designated by a number greater than $2m$; for otherwise the joining line would necessarily cut the axis. Thus, if 1 is supposed to be joined to the point $n$, we shall have $n<2m$. For a similar reason, if 2 is taken to be joined to $n'$, we shall have $n'<n$; for otherwise the branch $2\ldots n'$ would necessarily cut the branch $1\ldots n$. For the same reason, the point 3 will be joined to one of the points lying between 4 and $n'$. It is clear that if 3, 4, 5, etc. are supposed to be joined to $n''$, $n'''$, $n''''$, etc., then $n'''$ lies between 5 and $n''$, $n''''$ between 6 and $n'''$, etc. It is therefore evident that we shall eventually reach some point $h$ joined to the point $h+2$. Then the branch entering the circle at the point $h+1$ will necessarily cut the branch joining the points $h$ and $h+2$. Since one of these two branches belongs to the first line and the other to the second, it is now clear that the supposition is contradictory, and hence that an intersection of the first line with the second line must occur somewhere.
\SourceLeaf{034}
Combining this with what precedes, we conclude from all the investigations set out above that the theorem \textit{that every polynomial in one variable can be resolved into real factors of the first or second degree} has been proved with full rigour.

\Article{23} It is also not difficult to deduce from the same principles that there is not merely one intersection of the first line with the second, but at least $m$. It may, however, happen that the first line is cut by several branches of the second line at the same point, in which case the polynomial $X$ will have several equal factors. But since it is enough here to have proved the necessity of one intersection, I shall not dwell on this further, for brevity. For the same reason, I shall not pursue other properties of these lines in more detail here: for example, that their intersections always occur at right angles; or that, if several arms of both curves meet at the same point, there will be as many arms of the first line as of the second, arranged alternately and intersecting at equal angles, etc.

Finally, I observe that it is by no means impossible to present the preceding proof, which I have here built on geometric principles, in a purely analytic form as well. But I thought that the presentation set out here would be less abstract, and that the true force of the proof is brought much more clearly before the eyes here than could be expected from an analytic proof.

To finish, I shall also indicate another method of proving our theorem. At first sight it will seem very different, not only from the preceding proof, but also from all the other proofs described above; yet, if one looks at its essence, it is in fact the same as \textsc{d'Alembert}'s. I leave it to the experts to compare it with his and to examine the parallel between the two; it is added solely for their sake.

\Article{24} Above the plane of Fig. 4, I suppose the first and second surfaces to be drawn in the same way as above, relative to the axis $CG$ and the fixed point $C$. Take any point on some branch of the first line, that is, where $T=0$ (for example, any point $M$ lying on the axis). Unless $U=0$ holds there too, proceed from this point along the first line in the direction in which the absolute\SourceLeaf{035}
 value of $U$ decreases. If at the point $M$ the absolute value of $U$ happens to decrease in both directions, it is immaterial which way you proceed; I shall show at once what is to be done if $U$ increases in both directions. It is therefore clear that, as you continue along the first line, you must eventually reach either a point where $U=0$, or one where the value of $U$ becomes a minimum, for example the point $N$. In the former case, what was sought has been found. In the latter, it can be proved that several branches of the first line intersect at this point (an even number of branches, in fact), half of which are such that, if you turn into one of them (in either direction), the value of $U$ continues to decrease. (For brevity, I must omit the proof of this theorem, which is more lengthy than difficult.) You can therefore proceed again along this branch until $U$ either becomes $=0$ (as happens at $P$ in Fig. 4) or again becomes a minimum. Then you will turn aside again, and must eventually reach a point where $U=0$.

\Needspace{13\baselineskip}
One might raise against this proof the doubt whether it is possible that, however far you proceed and although the value of $U$ always decreases, these decreases nevertheless become continually slower, and that value never reaches some limit. This objection would correspond to the fourth in Article 6. But it would not be difficult to assign a boundary beyond which the value of $U$ must not only change ever more rapidly, but can no longer \textit{decrease}, so that the value 0 must necessarily already have occurred before you reach this boundary. I reserve for another occasion a fuller treatment of this and the remaining matters that I have only been able to indicate in this proof.

\begin{center}
\rule{8cm}{0.3pt}\\[8pt]
\textit{We discovered the principles on which this proof rests at the beginning of October} 1797.\\[8pt]
\rule{5cm}{0.3pt}
\end{center}
